\chapter{IEC 61131-3 em profundidade: POUs, tipagem e tarefas}
\label{cap:iec-modelo}

\begin{objetivos}
Ao final deste capítulo você deve ser capaz de:
\begin{itemize}
  \item descrever o modelo de software da norma, da configuração ao POU;
  \item diferenciar função, bloco de função e programa;
  \item configurar tarefas contínua, periódica e por evento, e explicar o papel do \emph{watchdog};
  \item classificar variáveis quanto ao escopo e à ligação com o hardware;
  \item usar os atributos de variável, em especial a retenção;
  \item dizer o que a norma \emph{não} padroniza --- e por que isso é a origem dos problemas de
        portabilidade.
\end{itemize}
\end{objetivos}

\section{O modelo de software}

A norma não padroniza apenas a forma das linguagens: padroniza também como o software se organiza.
Essa organização é hierárquica e é o que permite que um programa grande seja dividido em partes
gerenciáveis.

\begin{figure}[H]
  \centering
  \includegraphics[width=0.80\linewidth]{11-modelo-software.pdf}
  \caption{Modelo de software da norma: configuração, recursos, tarefas e POUs.}
  \label{fig:modelo-software}
\end{figure}

\paragraph{Configuração.} Corresponde ao software necessário a um controlador ou a um conjunto de
controladores para cumprir suas funções. Define os elementos que interagem entre si, com suas
configurações individuais, e a configuração total resultante.

\paragraph{Recursos.} Qualquer elemento com capacidade de processamento dentro da configuração,
capaz de executar programas. Pode ser físico --- a CPU do processador, uma interface de operação, um
\emph{gateway} de comunicação --- ou virtual, quando há compartilhamento de memória de processamento
por softwares distintos.

\paragraph{Tarefas.} Controlam a execução dos programas, de forma periódica ou por disparo de evento.
É a tarefa que define \emph{quando} cada parte do programa roda (Seção~\ref{sec:tarefas}).

\paragraph{POUs.} A \emph{unidade de organização de programa} é a forma definida pela norma para
implementar o software: associação de variáveis e instruções, em qualquer das linguagens. O software
aplicativo do controlador consiste na criação e associação de POUs entre si --- POUs podem ser
formadas a partir de outras POUs, o que permite estruturar o programa e reutilizar código.

\index{modelo de software}\index{POU}

\section{Função, bloco de função e programa}

\begin{table}[H]
  \centering
  \caption{As três unidades de organização de programa.}
  \label{tab:pous}
  \begin{tabular}{p{2.6cm} p{2.4cm} p{2.4cm} p{7.0cm}}
    \toprule
    POU & Memória interna & Saídas & Exemplos e uso \\
    \midrule
    Função (FUN)
      & Não tem: gera o resultado a cada execução, a partir das entradas
      & Um único valor
      & \texttt{ADD}, \texttt{SQRT}, \texttt{SIN}, \texttt{COS}, \texttt{GT}, \texttt{MIN},
        \texttt{MAX}, \texttt{AND}, \texttt{OR}; também funções definidas pelo usuário.
        Uso geral e sem estado \\
    Bloco de função (FB)
      & Tem: os dados persistem entre execuções --- é isso que o diferencia da função
      & Um ou mais valores
      & PID, temporizadores, contadores e blocos criados para funções específicas, como controle de
        motores. Podem compor blocos mais complexos \\
    Programa (PRG)
      & Tem estado próprio
      & ---
      & Construído a partir de funções e blocos de função, em qualquer linguagem da norma; pode
        acessar diretamente entradas e saídas e comunicar-se com outros programas. Suas partes podem
        ser controladas por tarefas diferentes \\
    \bottomrule
  \end{tabular}
\end{table}

A distinção função \emph{versus} bloco de função é a que mais causa erro em prova e em projeto: a
função é um cálculo; o bloco de função é um \emph{estado que evolui}. Um temporizador precisa
lembrar quanto tempo já acumulou --- por isso é bloco de função, e não função.

\begin{atencao}
Se você declarar o mesmo bloco de função em dois pontos do programa esperando compartilhar o estado,
o resultado não será o esperado: cada \emph{instância} tem seu próprio estado. Compartilhar estado se
faz por uma única instância declarada e usada em todo o programa. Este é um dos defeitos mais difíceis
de localizar em lógica de temporização.
\end{atencao}

\section{Tarefas}
\label{sec:tarefas}

\begin{table}[H]
  \centering
  \caption{Tipos de tarefa e consequência no comportamento do programa.}
  \label{tab:tipos-tarefa}
  \begin{tabular}{p{3.0cm} p{11.6cm}}
    \toprule
    Tarefa & Como se comporta \\
    \midrule
    Contínua
      & Executa em segundo plano, com menor prioridade: todo tempo de CPU não alocado a outras
        operações (comunicação, tarefas periódicas) é usado por ela. Ao terminar um ciclo de
        varredura, recomeça imediatamente. Um projeto não exige tarefa contínua; se houver, só pode
        existir uma \\
    Periódica
      & Executa em intervalos regulares de tempo. Quando o período se renova, interrompe qualquer
        tarefa de prioridade mais baixa e, ao terminar, devolve o controle ao ponto em que a tarefa
        anterior foi interrompida \\
    Por evento
      & Executa apenas quando ocorre o disparo configurado. Ao ocorrer, interrompe todas as tarefas
        de prioridade mais baixa, executa uma vez e devolve o controle ao ponto onde a anterior
        parou \\
    \bottomrule
  \end{tabular}
\end{table}

Cada tarefa tem um controle de tempo de execução configurável, o \emph{watchdog} (WDT). Ao estourar
o tempo de ciclo da CPU, o processamento para em condição de falha e o evento é sinalizado ao
operador de forma visual e, às vezes, sonora.

Separar o programa em tarefas é o principal instrumento de engenharia para conciliar exigências
diferentes no mesmo controlador: uma malha de controle que precisa rodar a cada 10\,ms não deve
esperar a rotina de comunicação, e um laço de diagnóstico lento não deve consumir o tempo do
controle.

\index{tarefa}\index{watchdog}

\section{Variáveis}

\subsection{Escopo}

\begin{itemize}
  \item \textbf{globais} --- declaradas acima de todos os programas; podem ser usadas pelo programa
        principal e por todas as rotinas e sub-rotinas;
  \item \textbf{locais} --- declaradas no início de uma rotina e visíveis apenas por ela.
\end{itemize}

A norma admite que uma variável global seja acessada por outra configuração ou mesmo por outro
controlador, desde que devidamente declarada em cada um.

\subsection{Ligação com o hardware}

\begin{itemize}
  \item \textbf{internas} --- sem vínculo direto com o hardware de E/S. Servem para resultados
        intermediários de cálculo, comparações e variáveis auxiliares. Não fazem parte da memória
        imagem de E/S;
  \item \textbf{externas} --- com vínculo direto com o hardware: são as variáveis de pontos de
        entrada e saída físicos, que recebem ou enviam valores para dispositivos do campo.
\end{itemize}

Confundir as duas é a origem de uma pergunta frequente em manutenção: ``por que este bit não aparece
no diagnóstico de E/S?''. Se ele não é externo, não tem ponto físico correspondente.

\subsection{Atributos de variável}

\begin{table}[H]
  \centering
  \caption{Atributos de variável previstos pela norma.}
  \label{tab:atributos-variavel}
  \begin{tabular}{p{2.4cm} p{12.2cm}}
    \toprule
    Atributo & Efeito \\
    \midrule
    \texttt{RETAIN}
      & Preserva o valor após o desligamento: é a variável retentiva do
        Capítulo~\ref{cap:memoria}. Use para contagens de turno, estados de processo que devem
        sobreviver à falta de energia \\
    \texttt{CONSTANT}
      & Valor fixo, não pode ser alterado em tempo de execução. Use para limites de projeto e
        parâmetros de calibração \\
    \texttt{AT}
      & Liga a variável a um endereço físico direto de E/S, dispensando o mapeamento por imagem.
        Recurso poderoso e traiçoeiro: o valor \emph{não} passa pela memória imagem \\
    \bottomrule
  \end{tabular}
\end{table}

E ainda os tipos de variável: global, direta (local), de mapeamento de E/S (entrada, saída, entrada
e saída), externa e temporária.

\begin{atencao}
Variável com \texttt{AT} não passa pela memória imagem: ela lê ou escreve o ponto físico no momento
do acesso, e não no início ou no fim do ciclo. Isso quebra o determinismo do programa e, em lógica
de segurança, é inaceitável. Use \texttt{AT} apenas com justificativa registrada.
\end{atencao}

\subsection{Tipagem e alocação}

Os tipos elementares e derivados foram apresentados no Capítulo~\ref{cap:memoria}. Além deles, a
norma permite criar novos tipos: dados estruturados que reúnem variáveis de tipos diferentes, tipos
derivados de um tipo já existente, tipos enumerados, subfaixas de valor com limites e arranjos.

Quanto à memória, os controladores atuais fazem \emph{alocação dinâmica} das variáveis: elas são
criadas com o tipo definido na declaração e não precisam ter tamanho fixo nem ocupar posições
sequenciais. Isso muda a pergunta de projeto: em vez de calcular endereços, declara-se a variável
com o tipo adequado e o ambiente resolve a alocação.

\section{Estruturação, modularização e reutilização}

O modelo de software é uma combinação de conceitos que dá estrutura aos projetos. Três princípios o
sustentam:

\begin{enumerate}
  \item \textbf{modularização} --- dividir o problema em partes com responsabilidade definida;
  \item \textbf{estruturação} --- organizar essas partes em hierarquia, com interfaces claras;
  \item \textbf{reutilização} --- escrever uma vez e usar em vários pontos, por instanciação de
        blocos de função.
\end{enumerate}

Trocando em miúdos: o programa deixa de ser uma sequência longa de linhas e passa a ser um conjunto
de blocos nomeados, cada um testável isoladamente. É o que torna possível dividir o trabalho entre
programadores e, depois, entregar a manutenção a outra equipe.

\section{O que a norma não define}

Este ponto é decisivo para o capítulo seguinte. A norma \emph{não} define:

\begin{itemize}
  \item o endereçamento físico das entradas e saídas (\texttt{\%I0.0}, \texttt{I:1/0},
        \texttt{DI1}\ldots);
  \item o ambiente de programação, sua interface e seu formato de arquivo;
  \item os protocolos de rede usados para programar e para comunicar;
  \item a implementação dos blocos de sistema e de comunicação específicos do fabricante.
\end{itemize}

Portanto, ``conforme IEC 61131-3'' descreve um conjunto de garantias sobre \emph{software aplicativo}
--- e nada além disso. Todo o resto continua sendo decisão do fabricante, e é exatamente aí que a
portabilidade se perde.

\begin{pratica}
Em um controlador com ambiente de programação disponível:
(a) crie uma função que converta uma leitura de 4 a 20\,mA em unidade de engenharia;
(b) crie um bloco de função que acumule o tempo em que uma entrada permanece acionada, com
\texttt{RETAIN};
(c) instancie o bloco em dois pontos do programa e verifique se os estados são independentes;
(d) crie duas tarefas --- uma periódica de 20\,ms para o controle e uma lenta para diagnóstico --- e
observe a diferença no comportamento do programa.
\end{pratica}

\begin{exercicios}
\begin{enumerate}[label=\alph*)]
  \item Desenhe a hierarquia do modelo de software e explique o papel de cada nível.
  \item Por que o temporizador é bloco de função e não função? Responda em termos de memória.
  \item Explique a diferença entre tarefa contínua, periódica e por evento, dando um uso real para
        cada uma.
  \item O que o \emph{watchdog} de uma tarefa protege, e o que o operador deve fazer quando ele
        dispara?
  \item Classifique como interna ou externa: (a) contador de peças do turno; (b) estado do
        pressostato de linha; (c) resultado intermediário de um cálculo de vazão; (d) comando da
        válvula de saída.
  \item Um técnico precisa que a contagem de ciclos não se perca na falta de energia. Que atributo de
        variável usar, e onde mais o mesmo atributo é necessário no sistema?
  \item Explique por que \texttt{AT} quebra o determinismo do ciclo de varredura.
  \item Liste três itens que a norma \emph{não} padroniza e explique como cada um deles pode impedir
        a troca de fabricante.
\end{enumerate}
\end{exercicios}
